% Phase-2 fix designs for defects R-011, R-012, R-013, R-015, R-016.
% Section 2 fragment of the gap-analysis document.  Plain LaTeX; compiles
% under amsmath/amssymb; each design fits one page.  Line anchors refer to
% the v1.0 manuscript source/formal-logic-reference.tex (14780 lines).

\subsection*{R-011 --- CNF half of the normal-form theorem (v1.0 lines 1919--1966; defective half 1958--1965)}

\paragraph{Defect.} The CNF half of the theorem at 1919--1923 reads ``Apply
the disjunctive construction to \(\neg\varphi\), negate the result, and use
De Morgan equivalences'' (1958--1959).  Against the string-valued CNF
definition at 1878--1897 (a CNF formula is \(\top\), or the right-associated
conjunction \(N_1\) of clause interpretations \(\operatorname{form}(C_i)\),
each clause listed in strictly increasing code order), this fails: for valid
\(\varphi\) the DNF of \(\neg\varphi\) is \(\bot\), and \(\neg\bot\) is not a
CNF string; for variable-free false \(\varphi\) it yields \(\neg\top\)
instead of \(\bot=\operatorname{form}(\varnothing)\); and ``use De Morgan''
silently performs double-negation elimination and re-association that the
definition does not supply.

\paragraph{Repair.} Replace 1958--1965 by a dual direct construction
mirroring 1924--1956, plus one auxiliary lemma inserted after 1897.
\emph{Auxiliary lemma} (complete proof; two inductions on the definitional
recursions at 1855--1870 and 1883--1890): (i) a valuation \(w\) satisfies
\(\operatorname{form}(C)\) iff \(w\) satisfies some literal in \(C\)
(induction on \(|C|\) along \(D_k=L_k\), \(D_i=(L_i\lor D_{i+1})\); the case
\(C=\varnothing\) is the falsity of \(\bot\)); (ii) \(w\) satisfies \(N_1\)
iff \(w\) satisfies \(\operatorname{form}(C_j)\) for every \(j\le m\)
(induction on \(m\) along the \(N_i\) recursion).
\emph{Construction.}  Let \(p_1,\ldots,p_n\) be the variables of
\(\varphi\) in increasing index order; reuse the DNF half's valuations on
this finite set and the coincidence lemma (1925--1930).  For a valuation
\(v\) on \(\{p_1,\ldots,p_n\}\), let \(\bar\ell^{\,v}_i\) be \(\neg p_i\) if
\(v(p_i)=\mathsf{T}\) and \(p_i\) if \(v(p_i)=\mathsf{F}\), and put
\(K_v:=\{\bar\ell^{\,v}_1,\ldots,\bar\ell^{\,v}_n\}\).  The coding at
1849--1854 gives the variable of index \(i\) the code \(2i\) and its negation
the code \(2i+1\) --- both strictly increasing in \(i\) --- and the
enumeration \(p_1,\ldots,p_n\) is in increasing index order (the indices need
not form an initial segment), so any literal on \(p_j\) has smaller code than
any literal on \(p_k\) for \(j<k\): if their indices are \(i_j<i_k\), then
\(2i_j+1<2i_k\).  The definitional listing of \(K_v\) in increasing code
order is therefore exactly \(\bar\ell^{\,v}_1,\ldots,\bar\ell^{\,v}_n\):
\(\gamma_v:=\operatorname{form}(K_v)\) needs no re-ordering.  By lemma (i),
\(w\) falsifies \(\gamma_v\) iff \(w(p_i)=v(p_i)\) for all \(i\le n\).  Let
\(W_\varphi\) be the set of valuations on \(\{p_1,\ldots,p_n\}\) falsifying
\(\varphi\), enumerated lexicographically (\(\mathsf{F}<\mathsf{T}\), as at
1945--1948) as \(v_1,\ldots,v_m\).  Define \(C_\varphi:=\top\) if \(m=0\),
and otherwise \(C_\varphi:=N_1\) built from the clauses \(C_j:=K_{v_j}\) by
the recursion at 1883--1890.  Then \(C_\varphi\) is literally a CNF string
of 1878--1897, including the degenerate cases: valid \(\varphi\) gives
\(m=0\), the canonical \(\top\); variable-free true \(\varphi\) gives
\(n=m=0\), again \(\top\); variable-free false \(\varphi\) gives \(n=0\),
\(m=1\), \(K_{v_1}=\varnothing\), so
\(C_\varphi=\operatorname{form}(\varnothing)=\bot\), the CNF whose single
clause is empty.  \emph{Equivalence:} for a total valuation \(w\) with
restriction \(w_0\) to \(\{p_1,\ldots,p_n\}\) and \(m>0\): by lemma (ii)
and (i), \(w\) satisfies \(C_\varphi\) iff for every \(j\le m\) it satisfies
some literal of \(K_{v_j}\), iff \(w_0\ne v_j\) for all \(j\), iff
\(w_0\notin W_\varphi\), iff \(w_0\) satisfies \(\varphi\), iff
(coincidence) \(w\) satisfies \(\varphi\).  If \(m=0\), every valuation on
the variable set satisfies \(\varphi\), so by coincidence every \(w\)
satisfies both \(\varphi\) and \(\top\).  The theorem keeps its plain proof
environment; the DNF half is unchanged except that it may cite lemma (i).

\paragraph{Acceptance.} (1) For the three degenerate inputs (valid;
variable-free true; variable-free false) the constructed strings are
\(\top\), \(\top\), \(\operatorname{form}(\varnothing)\), each satisfying
the definition at 1878--1897 verbatim.  (2) The CNF half invokes only the
definitional recursions, the auxiliary lemma, and the coincidence lemma ---
no De Morgan, double negation, idempotence, or associativity.  (3) The
auxiliary lemma carries a complete proof by induction on both recursions.

\subsection*{R-012 --- Mix rule, measure, and honest corollary statuses (v1.0 lines 9505--9548, 9555--9565, 9575--9586)}

\paragraph{Defect.} Step 3 of the cut-elimination sketch (9528--9532)
appeals to ``the multicut (mix) rule'' and to elimination ``by the same
lexicographic measure,'' but no mix rule is displayed anywhere, the measure
at 9509--9511 is defined only for cut, and the equipollence of the cut and
mix calculi is unstated.  The corollaries at 9555 and 9575 carry unlabeled
proof environments, implying complete proofs, yet both consume the sketched
theorem, contravening the status policy at 313--319.

\paragraph{Repair.} (1) \emph{Display mix} after 9511, in the notation of
9200--9224, where \(A^{k}\) abbreviates \(k\) comma-separated occurrences
of \(A\).  For \(\mathsf{LK}\), with \(m,n\ge 1\) and \(\Delta,\Sigma\)
permitted to contain further occurrences of \(A\):
\[
  \frac{\Gamma\Rightarrow\Delta,A^{m}\qquad A^{n},\Pi\Rightarrow\Sigma}
       {\Gamma,\Pi\Rightarrow\Delta,\Sigma}\;\textsc{Mix}.
\]
For \(\mathsf{LJ}\) (single succedent, so the left premise's succedent
multiplicity \(m\) is fixed at one): from \(\Gamma\Rightarrow A\) and
\(A^{n},\Pi\Rightarrow C\) with \(n\ge1\),
conclude \(\Gamma,\Pi\Rightarrow C\).
(2) \emph{Equipollence lemma} (new; status: complete proof).  Cut is the
instance \(m=n=1\).  Conversely, mix is derived from the displayed
primitives: \(m-1\) applications of \textsc{CR} to the left premise,
\(n-1\) of \textsc{CL} to the right, then one cut.  Hence the cut and mix
formulations derive the same sequents (\(\mathsf{LJ}\): \textsc{CL} only),
so eliminating mix eliminates cut.
(3) \emph{Restate the measure} (replace 9505--9511): a mix instance with
mix formula \(A\) and immediate subderivations \(D_1,D_2\) has rank
\(r(A)\) and height \(h(D_1)+h(D_2)\); the reduction is controlled by the
lexicographic order on \(\langle r(A),\,h(D_1)+h(D_2)\rangle\), with cut
the case \(m=n=1\).
(4) \emph{Preface the sketch} (after 9519) with: ``By the equipollence
lemma, first replace every cut in the given derivation by the mix instance
\(m=n=1\); the reduction below eliminates mixes, and steps 1, 2, 5, and 6
apply verbatim with `mix' for `cut' and the mix formula for the cut
formula.''  \emph{Rewrite sketch step 3} (9528--9532): if a premise of a mix ends
in a contraction of the mix formula (\textsc{CL} above the right premise,
\textsc{CR} above the left), replace that premise by the contraction's
premise and raise the corresponding multiplicity by one: rank unchanged,
height decreased by one.  Add to step 4: when the mix formula is principal
on both sides with multiplicities \((m,n)\), the principal reduction yields
mixes on proper subformulas (smaller rank) plus residual mixes on remaining
occurrences at equal rank but smaller height; the full case table remains
the omitted verification (citation at 9548 unchanged).
(5) \emph{Honest statuses.}  Append \textit{Status: proof sketch.}\ to the
theorem statement (9513--9517).  Split the corollary at 9555: a new lemma
``every cut-free propositional \(\mathsf{LK}\)- or
\(\mathsf{LJ}\)-derivation has the subformula property'' takes the existing
induction (9560--9565) as a complete proof; the corollary (existence of a
cut-free derivation with the property) is restated with \textit{Status:
proof sketch.}, its proof converted to the sketch environment with the
dependency sentence modeled on 7728--7729: ``This derivation uses the
preceding proof sketch of the cut-elimination theorem; a complete proof
supplies its reduction case table.''  The corollary at 9575 keeps its
argument verbatim, converted to the sketch environment with the same
dependency sentence and \textit{Status: proof sketch.}\ on its statement.
When cut elimination is proved in full, exactly these two status lines and
two sentences flip to complete proof.

\paragraph{Acceptance.} (1) Mix is displayed for both calculi with explicit
multiplicities \(m,n\ge1\), and the equipollence lemma is proved using only
\textsc{CL}, \textsc{CR}, and \textsc{Cut} as displayed at 9204--9224 and
9307--9375.  (2) The sketch opens by converting every cut to a mix via the
equipollence lemma, so no step manipulates an object the measure does not
cover; step 3 names only displayed rules and exhibits the measure decrease
(rank equal, height \(-1\)).  (3) The theorem and both
corollaries carry \textit{Status: proof sketch.}; the cut-free-subformula
lemma and the equipollence lemma carry \textit{Status: complete proof.}; no
complete-proof item in the chapter cites the cut-elimination theorem.

\subsection*{R-013 --- Repetition and vacuous discharge in \(\mathsf{ND}_{\mathsf{CPL}}\) (v1.0 lines 2006--2021, 2351--2379, 2443--2455; first-order model 3799--3812)}

\paragraph{Defect.} Line 2020--2021 licenses, in prose only, ``Repetition
of an already available formula is permitted,'' though the calculus is
defined (2006--2018) as a tree generated by the displayed rules and no
repetition rule is displayed; and the definition never says whether a
discharging rule may discharge zero occurrences.  Clause 2 of the closure
lemma (2360--2362) and the completeness theorem (2446--2449) apply
\(\neg I\)/\(\mathsf{RAA}\) to derivations that may contain no open
occurrence of the discharged assumption, and clause 5 (2374--2375) says
``assume \(\varphi\), repeat \(\psi\), and apply \(\to I\).''

\paragraph{Repair.} Eliminate the uses, matching 3799--3812.
(a) Delete the repetition sentence (2020--2021).  Append to the derivation
definition (after 2011): ``An application of \(\to I\), \(\neg I\), or
\(\mathsf{RAA}\) carries a label \(i\) and discharges exactly the open
occurrences marked \(i\) of the bracketed assumption above its premise; an
application of \(\lor E\) carries two labels \(i,j\), as displayed, and
discharges exactly the open occurrences of the first disjunct marked \(i\)
above its second premise and the open occurrences of the second disjunct
marked \(j\) above its third premise.  Every label carried by such an
application discharges at least one occurrence.  Marking any set of open
occurrences of one formula with a common label fresh from every label in
the derivation preserves the derivation and the occurrences discharged by
each of its rule applications, since no rule application carries the fresh
label; this merging relabeling is distinct from the injective renamings
above, and is the operation used when a newly added root inference is to
discharge several open occurrences at once.  There is no repetition step:
where an informal presentation repeats a formula, the tree reuses the node
already present.''
(b) Closure clause 2 (2358--2362), mirroring 3803--3807: if
\(\varphi\notin\Delta\), maximality yields a derivation \(D\) of \(\bot\)
with open assumptions in \(\Delta\cup\{\varphi\}\); because \(\Delta\) is
consistent, at least one open occurrence of \(\varphi\) remains in \(D\);
mark all open occurrences of \(\varphi\) in \(D\) with a common fresh label
\(i\) (the merging relabeling of (a)), and apply \(\neg I\) with label
\(i\) at the root, discharging every such occurrence, to obtain
\(\Delta\vdash\neg\varphi\); clause 1 gives \(\neg\varphi\in\Delta\).
(c) Closure clause 5 (2374--2375), replacing repetition by a conjunction
detour: if \(\psi\in\Delta\), assume \(\varphi\) with label \(i\) and take
\(\psi\) as an open assumption; \(\land I\) gives \(\varphi\land\psi\);
\(\land E_2\) gives \(\psi\); \(\to I\) discharges the single open
occurrence marked \(i\), concluding \(\varphi\to\psi\) with open
assumptions in \(\Delta\).  The discharge is of an occurrence genuinely
present, and only displayed rules occur.
(d) Completeness (2446--2449): suppose a derivation \(D\) of \(\bot\) has
open assumptions in \(\Gamma\cup\{\neg\varphi\}\).  If no open occurrence
of \(\neg\varphi\) appeared in \(D\), then \(D\) would witness
\(\Gamma\vdash\bot\), and \(\bot E\) would give \(\Gamma\vdash\varphi\),
contrary to the supposition \(\Gamma\nvdash\varphi\).  So at least one open
occurrence remains; mark all open occurrences of \(\neg\varphi\) with a
common fresh label \(i\) (the merging relabeling of (a)) and apply
\(\mathsf{RAA}\) with label \(i\) at the root, discharging every such
occurrence, to obtain \(\Gamma\vdash\varphi\) --- again contrary to
the supposition.  Hence \(\Gamma\cup\{\neg\varphi\}\) is consistent.
(e) Sweep every displayed \(\lor E\)/\(\to I\)/\(\neg I\)/\(\mathsf{RAA}\)
application in the ND chapters, examples, exercises, and solutions for
compliance (the examples at 2077--2110 already comply; the intuitionistic
calculus, defined at 2106--2109 by omitting \(\mathsf{RAA}\), inherits the
convention; 3799--3812 stands unchanged and now matches, its ``discharging
every such occurrence'' being an instance of the merging relabeling
followed by a single labeled application).

\paragraph{Acceptance.} (1) ``Repetition'' no longer occurs in the
\(\mathsf{ND}_{\mathsf{CPL}}\) section, and the derivation definition
states explicitly both the at-least-one-occurrence discharge convention
(with \(\lor E\) carrying two labels) and the fresh-label merging
relabeling used by the closure and completeness proofs.
(2) The proofs of closure clauses 2 and 5 and of completeness contain,
respectively, the non-vacuity argument and the conjunction detour, and
invoke no rule outside the display at 2022--2065.  (3) An audit of all
discharging-rule applications in displayed derivations finds none that
discharges zero occurrences.

\subsection*{R-015 --- Reduction of arbitrary embeddings to substructure inclusions (v1.0 lines 5989--5992, 6009--6012; supporting definitions 5042--5057, 5093--5103)}

\paragraph{Defect.} Model completeness is defined at 5989 as ``every
embedding between its models is elementary,'' but the \(\mathsf{DLO}\)
corollary's sketch verifies only inclusions: ``If
\(\mathcal M\subseteq\mathcal N\) are models of \(\mathsf{DLO}\)
\(\ldots\) This is model completeness'' (6009--6012).  The passage from
inclusions to arbitrary embeddings is missing, and the manuscript nowhere
proves that satisfaction is isomorphism-invariant or that an embedding is
an isomorphism onto its image.

\paragraph{Repair.} Insert Lemmas 1--2 after 5125 (end of the
Substructures section) and Lemma 3 after 5992; patch 6011--6012.
\emph{Lemma 1 (Isomorphism invariance; complete proof).}  If
\(h:\mathcal M\to\mathcal N\) is an isomorphism, then for every formula
\(\varphi(x_1,\ldots,x_k)\) and all \(a_1,\ldots,a_k\in M\):
\(\mathcal M\models\varphi(\bar a)\iff\mathcal N\models\varphi(h\bar a)\);
in particular isomorphic structures satisfy the same sentences.
\emph{Route:} induction on \(\varphi\); atomic case by the
term-preservation lemma (5060--5076), injectivity for equality atoms, and
preservation with reflection (5042--5052) for relation atoms; Boolean cases
by induction hypothesis; for \(\exists y\), a witness \(a'\in M\) maps to
\(h(a')\), and conversely a witness \(b\in N\) is \(h(a')\) by
surjectivity; \(\forall\) dually.
\emph{Lemma 2 (Image of an embedding; complete proof).}  If
\(h:\mathcal M\to\mathcal N\) is an embedding, then \(h(M)\) contains
\(c^{\mathcal N}\) for every constant and is closed under every
\(f^{\mathcal N}\); the structure \(h[\mathcal M]\) with domain \(h(M)\),
constants \(c^{\mathcal N}\), the restrictions of the \(f^{\mathcal N}\),
and relations \(R^{\mathcal N}\) restricted to \(h(M)\) is a substructure
of \(\mathcal N\) per 5093--5103; and \(h\) with codomain \(h(M)\) is an
isomorphism \(\mathcal M\to h[\mathcal M]\).  \emph{Route:}
\(c^{\mathcal N}=h(c^{\mathcal M})\in h(M)\) and
\(f^{\mathcal N}(h\bar a)=h(f^{\mathcal M}(\bar a))\in h(M)\) by the
homomorphism clauses; the inclusion preserves and reflects atomic formulas
because interpretations are restrictions, so \(h[\mathcal M]\subseteq
\mathcal N\); the corestricted \(h\) is injective (given), surjective by
construction, preserves and reflects relations since
\(R^{h[\mathcal M]}=R^{\mathcal N}\cap h(M)^{k}\) while \(h\) reflects
\(R^{\mathcal N}\); hence a surjective embedding, an isomorphism per
5055--5057.
\emph{Lemma 3 (Reduction to inclusions; complete proof).}  If
\(\mathcal M\preccurlyeq\mathcal N\) whenever
\(\mathcal M\subseteq\mathcal N\) are both models of \(T\), then \(T\) is
model complete.  \emph{Route:} an embedding \(h:\mathcal M\to\mathcal N\)
between models of \(T\) factors (Lemma 2) as the isomorphism
\(h_0:\mathcal M\to h[\mathcal M]\) followed by the inclusion
\(h[\mathcal M]\subseteq\mathcal N\); Lemma 1 gives
\(h[\mathcal M]\models T\); the hypothesis gives
\(h[\mathcal M]\preccurlyeq\mathcal N\); then for all \(\varphi,\bar a\):
\(\mathcal M\models\varphi(\bar a)\iff h[\mathcal M]\models
\varphi(h\bar a)\) (Lemma 1) \(\iff\mathcal N\models\varphi(h\bar a)\)
(elementary inclusion).  Hence \(h\) is elementary.
\emph{Patch 6011--6012:} replace ``This is model completeness.''\ by ``By
the reduction lemma, every embedding between models of \(\mathsf{DLO}\) is
elementary; this is model completeness.''  The corollary keeps
\textit{Status: proof sketch.}\ (its remaining omission is the literal-case
verification named at 6013--6014, not the embedding reduction).

\paragraph{Acceptance.} (1) The three lemmas appear with complete proofs
using only 5022--5057, 5093--5103, and the term-preservation lemma; Lemma
1's quantifier case invokes surjectivity explicitly, and Lemma 2 verifies
closure of \(h(M)\) before forming \(h[\mathcal M]\).  (2) The
\(\mathsf{DLO}\) sketch derives model completeness by citing Lemma 3.
(3) No passage concludes model completeness from inclusions alone.

\subsection*{R-016 --- Extension clause in the back-and-forth test (v1.0 lines 6068--6093; corollary 6097--6110)}

\paragraph{Defect.} The back-and-forth theorem (6068--6081) concludes only
``\(\mathcal M\cong\mathcal N\),'' but the corollary's proof asserts ``The
resulting isomorphism extends the original map'' (6108--6109) --- a
property of the proof's construction, which begins at an arbitrary
\(f_0\in\mathcal F\), not of the stated theorem.  The corollary's use is
therefore unlicensed.

\paragraph{Repair.} \emph{Restated theorem (6068--6081).}  Let
\(\mathcal M,\mathcal N\) be countable structures and \(\mathcal F\) a
nonempty set of finite partial isomorphisms from \(\mathcal M\) to
\(\mathcal N\) such that for every \(f\in\mathcal F\): (1) for every
\(a\in M\), some \(g\in\mathcal F\) extends \(f\) with \(a\) in its domain;
(2) for every \(b\in N\), some \(g\in\mathcal F\) extends \(f\) with \(b\)
in its range.  Then for every \(f_0\in\mathcal F\) there is an isomorphism
\(h:\mathcal M\to\mathcal N\) with \(f_0\subseteq h\) (extension as sets of
ordered pairs); in particular \(\mathcal M\cong\mathcal N\).
\emph{Proof edits (6083--6093):} replace ``begin with any
\(f_0\in\mathcal F\)'' by ``let \(f_0\) be the given member of
\(\mathcal F\)''; the recursion and union argument are unchanged; append
``Since \(f_0\subseteq f_1\subseteq\cdots\) and \(h=\bigcup_k f_k\), the
isomorphism extends \(f_0\).''  \textit{Status: complete proof}\ retained.
\emph{Corollary patch (6104--6110).}  Replace the final sentence by an
explicit application: the set \(\mathcal F\) of finite elementary maps from
\(\mathcal M\) to itself extending the given map \(f\) is nonempty
(\(f\in\mathcal F\)); every member is a finite partial isomorphism, since
elementary maps preserve and reflect atomic formulas and reflect equality,
hence are injective; clause (1) holds within \(\mathcal F\) because the
saturation-homogeneity theorem (5567--5571) extends any \(g\in\mathcal F\)
by any prescribed domain element to an elementary map still extending
\(f\); clause (2) holds by applying the same theorem to \(g^{-1}\) (the
inverse of a finite elementary self-map is elementary) and inverting the
result, which extends \(g\supseteq f\) and has the prescribed element in
its range.  The restated theorem applied with \(f_0:=f\) yields an
isomorphism \(h:\mathcal M\to\mathcal M\) with \(f\subseteq h\): an
automorphism extending \(f\).  \textit{Status: complete proof}\ retained.
The forward reference at 5563--5564 remains accurate.

\paragraph{Acceptance.} (1) The theorem statement contains the quantified
extension clause ``for every \(f_0\in\mathcal F\) there is an isomorphism
\(h\supseteq f_0\).''  (2) Its proof starts the recursion at the given
\(f_0\) and derives \(f_0\subseteq h\) from the chain union.  (3) The
corollary invokes the extension clause with \(f_0:=\) the given map, and
both results retain \textit{Status: complete proof.}
